\section{Related work}

\paragraph{Rain modeling} In their series of influential papers, Garg and Nayar provided a comprehensive overview of the appearance models required for understanding~\cite{garg2007vision} and synthesizing~\cite{garg2006photorealistic} realistic rain. In particular, they propose an image-based rain streak database~\cite{garg2006photorealistic} modeling the drop oscillations, which we exploit in our physics-based rendering framework. 
Other streak appearance models were proposed in~\cite{weber2015multiscale,barnum2010analysis} using a frequency model. Realistic rendering was also obtained with ray-tracing~\cite{rousseau2006realistic} or artistic-based techniques~\cite{tatarchuk2006artist,creus2013r4} but on synthetic data as they require complete 3D knowledge of the scene including accurate light estimation. Numerous works also studied generation of raindrops on-screen, with 3D modeling and ray-casting~\cite{roser2009video,roser2010realistic,halimeh2009raindrop,hao2019learning} or normal maps~\cite{porav2019can} some also accounting for focus blur.

\paragraph{Rain removal} Due to the problems it creates on computer vision algorithms, rain removal in images got a lot of attention initially focusing on photometric models~\cite{garg2004detection}. For this, several techniques have been proposed, ranging from frequency space analysis~\cite{barnum2010analysis} to deep networks~\cite{yang2017deep}. Sparse coding and layers priors were also important axes of research~\cite{li2016rain,luo2015removing,chen2013generalized} due to their facility to encode streak patches. Recently, dual residual networks have been employed~\cite{liu2019dual}.
Alternatively, camera parameters~\cite{garg2005does} or programmable light sources~\cite{de2012fast} can also be adjusted to limit the impact of rain on the image formation process. 
Additional proposals were made for the specific task of raindrops removal on windows~\cite{eigen2013restoring} or windshields~\cite{halimeh2009raindrop,porav2019can,hao2019learning}.

\paragraph{Unpaired image translation} An interesting solution to weather augmentation is the use of data-driven unpaired image translation frameworks. Zhang et al.~\cite{zhang2019image} proposed to use conditional GANs for rain removal. By proposing to add a cyclic loss to the learning process, CycleGAN~\cite{zhu2017unpaired} became a significant paper for unpaired image translation; they produced interesting results in weather and season translation. DualGAN~\cite{yi2017dualgan} uses similar ideas with differences in the network models. The UNIT~\cite{liu2017unsupervised}, MUNIT~\cite{huang2018multimodal}, and FUNIT~\cite{liu2019few} frameworks all, in one way or another, propose to perform image translation with the common idea that data from different sets have a shared latent space. They showed interesting results on adding and removing weather effect to images. 
Since the information in clear and rainy images is symmetrical, many unsupervised image translation approaches could produce decent visual results. 
In \cite{pizzati2020model}, Pizzati et al. learn to disentangle the scene from lens occlusions such as raindrops, which improves both realism and physical accuracy of the translations.
Another strategy for better qualitative translations is to rely on semantic consistency~\cite{li2018semantic,tasar2020semi2i}. 

\paragraph{Weather databases}

In computer vision, few images databases have precise labeled weather information. 
Of note for mobile robotics, the BDD100K~\cite{yu2018bdd100k}, the Oxford dataset~\cite{RobotCarDatasetIJRR}, and Wilddash~\cite{zendel2018wilddash} provide data recorded in various weather conditions, including rain. 
Other stationary camera datasets such as AMOS~\cite{jacobs07amos}, the transient attributes dataset~\cite{Laffont14}, the Webcam Clip Art dataset~\cite{lalonde2009webcam}, or the WILD dataset~\cite{narasimhan2002all} are sparsely labeled with weather information. 
The relatively new nuScenes dataset~\cite{caesar2019nuscenes} have multiple labeled scenes containing rainy images, but variation in rain intensities are not indicated. Gruber et al.~\cite{gruber-3dv-19} recently released a dataset with dense depth labels under a variety of real weather conditions produced by a controlled weather chamber, which inherently limits the variety of scenes (limited to four common scenarios) in the dataset. Note that \cite{bijelic2019seeing} also announced---but at the time of writing, not yet fully available---a promising dataset including heavy snow and rain events.
Still, datasets with rainy data are too small to train algorithms and there exists no dataset with systematically recorded rainfall rates and object/scene labels. 
The closest systematic works in spirit~\cite{johnson2016driving,khan2019procsy} evaluated the effect of simulating weather on vision, but did so in purely virtual environments (GTA and CARLA, respectively) whereas we augment real images.
Of particular relevance to our work, Sakaridis et al.~\cite{SDV18} propose a framework for rendering fog into images from the Cityscapes~\cite{Cordts2016Cityscapes} dataset. Their approach assumes a homogeneous fog model, which is rendered from the depth estimated from stereo. Existing scene segmentation models and object detectors are then adapted to fog. 
In our work, we employ the similar idea of rendering realistic weather on top of existing images, but we focus on rain rather than fog. 
